\section{池化操作与文本分类}

\subsection{Max Pooling（最大池化）}

卷积操作为每个位置产生了特征向量，但对于文本分类任务，我们需要一个\textbf{固定长度}的句子表示。最常用的方法是 \textbf{Max Pooling}（最大池化）。

\begin{figure}[H]
\centering
\includegraphics[width=0.95\textwidth]{slides-images/slide-020.jpg}
\caption{Max Pooling 与分类：对每个滤波器的所有位置取最大值，得到固定长度的特征向量\protect\footnotemark}
\end{figure}
\footnotetext{Slides 第21页。}

\begin{importantbox}{Max Pooling 的直觉}
可以将每个滤波器理解为一个\textbf{特征检测器}。例如，某个滤波器可能学会检测``使用第一人称代词''（匹配 I、my、we、our 等词）。Max Pooling 的作用是：\textbf{只要这个特征在文本的任何位置被强烈激活，就认为该特征存在}。这类似于一个``是否存在某模式''的二值检测器。
\end{importantbox}

具体地，设滤波器 $j$ 在所有位置产生的特征图为 $\mathbf{c}_j = [c_{j,1}, c_{j,2}, \ldots, c_{j,n-h+1}]$，则：

$$\hat{c}_j = \max\{\mathbf{c}_j\}$$

对 $m$ 个滤波器分别做 max pooling，得到最终的句子表示向量：

$$\mathbf{z} = [\hat{c}_1, \hat{c}_2, \ldots, \hat{c}_m] \in \mathbb{R}^m$$

\subsection{Average Pooling（平均池化）}

另一种池化方式是 \textbf{Average Pooling}。如果将滤波器理解为度量文本某种属性的程度（如``正式程度''），那么平均值可能比最大值更合理——它反映了整体水平而非峰值。

\begin{warningbox}{Max Pooling 通常优于 Average Pooling}
在实践中，如果只选择一种池化方式，Max Pooling 通常表现更好。这是因为神经网络学到的特征更接近``模式检测器''而非``程度度量器''。当然，也可以同时使用两种池化方式并拼接结果，这有时能带来额外的性能提升。
\end{warningbox}

\subsection{其他池化变体}

\subsubsection{Local Max Pooling}

不对整个序列做全局 max pooling，而是在局部窗口内做 max pooling。这样可以保留特征出现的大致位置信息，并捕获同一特征在多个位置的出现。

\begin{figure}[H]
\centering
\includegraphics[width=0.95\textwidth]{slides-images/slide-015.jpg}
\caption{Local Max Pooling（stride=2）：对相邻两个位置做 max pooling，保留位置信息\protect\footnotemark}
\end{figure}
\footnotetext{Slides 第16页。}

\subsubsection{$k$-Max Pooling}

保留每个通道中最大的 $k$ 个值（而不是只保留最大的 1 个）。这样可以检测同一特征在文本中出现的次数——如果某个特征在多个位置被激活，$k$-max pooling 能够保留这些信息。

\subsection{Stride（步幅）}

卷积操作中，滤波器不一定要逐位移动。\textbf{Stride}（步幅）控制滤波器每次移动的距离。stride=2 意味着隔一个位置计算一次，输出长度减半。这在滤波器较大时尤其有用，可以减少冗余计算。

\subsection{Dilation（膨胀卷积）}

膨胀卷积不改变滤波器大小，而是在滤波器元素之间插入间隔。例如，dilation=2 的 trigram 滤波器会查看位置 $\{i, i+2, i+4\}$ 而不是 $\{i, i+1, i+2\}$。这使得滤波器能覆盖更大的上下文范围，而不增加参数量。

\begin{figure}[H]
\centering
\includegraphics[width=0.95\textwidth]{slides-images/slide-017.jpg}
\caption{膨胀卷积示意：通过在滤波器元素间插入间隔来扩大感受野\protect\footnotemark}
\end{figure}
\footnotetext{Slides 第18页。}

\begin{knowledgebox}{膨胀卷积的应用场景}
膨胀卷积在语音处理中比在自然语言处理中更常用。在语音信号中，远距离的时间步之间可能存在有意义的关联，膨胀卷积可以用较少的参数捕获这种远程依赖。在文本中，由于词本身已经是较高层次的单位，膨胀卷积的使用相对较少。
\end{knowledgebox}

\subsection{PyTorch 实现}

在 PyTorch 中，1D 卷积通过 \texttt{nn.Conv1d} 实现：

\begin{lstlisting}[caption={PyTorch 中 1D 卷积的基本用法},language=Python]
import torch.nn as nn

# in_channels: 词向量维度
# out_channels: 滤波器数量
# kernel_size: n-gram 大小
conv = nn.Conv1d(in_channels=300, out_channels=100,
                 kernel_size=3, padding=1)

# Max pooling over time
pool = nn.AdaptiveMaxPool1d(1)

# 输入: (batch_size, embedding_dim, seq_len)
# 输出: (batch_size, num_filters, 1)
\end{lstlisting}

\subsection{本章小结}

池化操作将变长的卷积特征图压缩为固定长度的表示向量。Max Pooling 是最常用的方式，它将每个滤波器视为特征检测器。Local max pooling、$k$-max pooling、stride 和 dilation 提供了更丰富的选择，可以根据具体任务灵活使用。

%% ========== Kim (2014) CNN ==========

